\documentstyle[%  


righttag%  


,11pt%  


,amssymb%  


,russian%               remove in Eng.  


,izvru%                 Set izven% in Eng.  


]{amsart}  


  


%       Math definitions  


 


\newcommand{\A}{{\cal A}}  


\newcommand{\re}{\operatorname{Re}}  


\newcommand{\const}{\operatorname{const}}  


\newcommand{\tts}[1]{} 


  


%\hoffset=-15mm%                  For Eng.  


  


\setcounter{page}{77}  


\god{1998}  


\nomer{2~(429)} %                Remove in Eng.  


%\nomer{Vol.\,42, No.\,2}%        Set in Eng.  


% \str{75--81}%                    Set in Eng.  


\udk{517.537.2}  


\author{\it М.Н.\,Шеремета, Я.Я.\,Притула}  


% NB:   ^^ remove in Eng. 


\title{Максимум модуля и максимальный член\\ 


одного класса рядов Дирихле}  


  


\date{}  


% \pagestyle{myheadings}%         Set in Eng.  


  


\pagestyle{plain} %              Remove in Eng.  


\begin{document}  


  


% \thispagestyle{plain}%          Set in Eng.  


  


\maketitle  


 


 


{\bf 1.} Пусть $0<A\le+\infty$, а $\Pi(A)$ --- класс функций, 


аналитических в 


$\{s:0<\re s<A\}$ и ограниченных в каждой полосе 


$\{s:0<\sigma_1\le\re s\le \sigma_2<a\}$. 


Изучению функций из класса $\Pi(\infty)$ посвящена глава 


{III} монографии [1]. Классу $\Pi(\infty)$ принадлежат, в частности, 


целые 


(абсолютно сходящиеся в $\Bbb C$) ряды Дирихле с неотрицательными, 


возрастающими к $\infty$ показателями $\lambda_n$ и комплексными 


коэффициентами 


$b_n$. Связь между максимумом модуля и максимальным членом 


целого ряда Дирихле исследована в [2]. 


 


Здесь рассмотрим ряды Дирихле, показатели которых не обязательно 


монотонны и могут принимать значения из $[-\infty,+\infty)$, но 


будем считать, что среди $\lambda_n$ имеется бесконечно много 


положительных. 


Считаем также, что $1>|b_n|\downarrow 0$ ($n\to\infty$). Через 


$\Pi_D(A)$ обозначим класс рядов Дирихле 


\begin{equation} 


F(s)=\sum^\infty_{n=1}b_n\exp(s\lambda_n),\quad 


s=\sigma+it, 


\end{equation} 


абсолютно сходящихся в $\{s:0\le\sigma<A\}$, коэффициенты и показатели 


которых удовлетворяют приведенным выше условиям. При 


этом, если $0<A<+\infty$, то считаем, что ряд 


\begin{equation} 


\sum^\infty_{n=1}|b_n|\exp(\sigma\lambda_n) 


\end{equation} 


сходится для любого $\sigma\in[0,A)$ и расходится для любого 


$\sigma>A$, 


а если $A=+\infty$, то считаем, что ряд (2) сходится для всех 


$\sigma\ge 0$. 


Другими словами, $A$ является абсциссой абсолютной сходимости 


ряда (1). Условимся также, что $|b_n|\exp(\sigma\lambda_n)=0$ для всех 


$\sigma\ge 0$, 


если $\lambda_n=-\infty$. Так как ряд (2) сходится при $\sigma=0$, то 


$$ 


\tau=\inf\biggl\{\lambda>0:\sum^\infty_{n=1}|b_n|^\lambda<\infty 


\biggr\}<\infty. 


$$ 


Положим 


$$ 


\gamma=\underset{n\to\infty}{\varlimsup} 


\frac{\lambda_n}{\ln(1/|b_n|)}. 


$$ 


Тогда $\gamma\ge 0$ и имеет место 


 


\begin{lem} Если $\tau<1$, то 


\begin{equation} 


\frac{1-\tau}{\gamma}\le A\le\frac{1}{\gamma}. 


\end{equation} 


\end{lem} 


 


\begin{pf} Пусть $\gamma>0$ и $\sigma>\frac{1}{\gamma}$. 


Тогда для любого 


$\varepsilon\in(0,\gamma-1/\sigma)$ существует возрастающая 


последовательность $(n_k)$ натуральных 


чисел такая, что 


$\lambda_{n_k}\ge(\gamma-\varepsilon)\ln\frac{1}{|b_{n_k}|}$, и, значит, 


$$ 


|b_{n_k}|\exp(\sigma\lambda_{n_k})\ge\exp 


\biggl\{-\ln\frac{1}{|b_{n_k}|}+\sigma(\gamma-\varepsilon) 


\ln\frac{1}{|b_{n_k}|}\biggr\}\ge 1, 


$$ 


т.е. для $\sigma>\frac{1}{\gamma}$ ряд (2) расходится. Поэтому имеет 


место второе из 


неравенств (3), которое очевидно, если $\gamma=0$. 


 


Пусть теперь $\gamma<+\infty$. Выберем $\Bar\gamma>\gamma$ и 


$\Bar\tau\in(\tau,1)$. Тогда для 


всех 


$n\ge n_0(\Bar\gamma)$ имеет место неравенство 


$\lambda_n\le\Bar\gamma\ln(1/|b_n|)$ и, значит, для 


всех $\sigma\in\Bigl[0,\frac{1-\Bar\tau}{\gamma}\Bigr)$ 


$$ 


\sum^\infty_{n=1}|b_n|e^{\sigma\lambda_n}\le 


\sum^\infty_{n=1}|b_n|\exp\{-\sigma\Bar\gamma\ln|b_n|\} 


\le\sum^\infty_{n=1}|b_n|^{\Bar\tau}<\infty. 


$$ 


 


Отсюда вытекает, что $A\ge\frac{1-\Bar\tau}{\gamma}$, и в силу 


произвольности $\Bar\tau$ и $\Bar\gamma$ 


имеет место первое из неравенств (3), которое очевидно, если 


$\gamma=+\infty$. 


\end{pf} 


 


Заметим, что если $\tau=0$, то $A=1/\gamma$, а это возможно, например, 


если $\sqrt[\uproot{2}n]{|b_n|}\to 0$ ($n\to\infty$), ибо тогда 


$\sum\limits^\infty_{n=1}|b_n|^\lambda<\infty$ для 


любого 


$\lambda>0$. 


 


Заметим также, что оценки (3) точны. Так, например, если 


$b_n=1/n^2$ 


и $\lambda=2\ln n$, то $\tau=1/2$, $\gamma=1$ и 


$(1-\tau)\gamma=1/2=A<1=1/\gamma$. 


А если $b_n=1/n^2$, $n_k=k!$, $\lambda_{n_k}=2\ln n_k$ и $\lambda_n=-n$ 


при $n\ne n_k$, то 


$\tau=1/2$, $\gamma=1$ и $(1-\tau)\gamma=1/2<A=1=1/\gamma$, ибо для всех 


$0\le\sigma<1$ 


имеем 


$$ 


\sum^\infty_{n=1}|b_n|e^{\sigma\lambda_n}= 


\sum^\infty_{k=1}n^{-2}_k\exp(2\sigma\ln n_k)+ 


\sum^\infty_{n=1,\,n\ne n_k}n^{-2}\exp(-\sigma n)\le 


\sum^\infty_{k=1}(k!)^{-2(1-\sigma)}+ 


\sum^\infty_{n=1}n^{-2}<+\infty. 


$$ 


 


Для $0\le \sigma<A$ и $F\in\Pi_D(A)$ положим 


$M(\sigma,F)=\sup\{|F(\sigma+it)|:t\in\Bbb R\}$, 


и пусть $\mu(\sigma,F)=\max\{|b_n|\exp(\sigma\lambda_n):n\ge 1\}$ --- 


максимальный 


член ряда (1). Из доказательства теоремы 1.5 в [2] следует 


неравенство $\mu(\sigma,F)\le M(\sigma,F)$, являющееся аналогом 


известного неравенства Коши. 


 


Целью данной статьи является оценка $M(\sigma,F)$ через $\mu(\sigma,F)$ 


сверху для $F\in\Pi_D(A)$ как в случае $0<A<+\infty$, так и в случае 


$A=+\infty$. 


 


\begin{lem} Пусть $b_n\in\Bbb C$, $\lambda_n\in\Bbb R$ $(n\ge 1)$, 


$p\ge 0$, $q>0$, а ряд Дирихле 


{\em (1)} абсолютно сходится в точках $s=\sigma$ и 


$s=\frac{1}{q}(\sigma+p)$. Тогда, если 


\begin{equation} 


\sum^\infty_{n=1}|b_n|^{1-q}\exp(-p\lambda_n)=K_0<\infty, 


\end{equation} 


то 


\begin{equation} 


M(\sigma,F)\le K_0\mu 


\biggl(\frac{\sigma+p}{q},\,F \biggr)^q. 


\end{equation} 


\end{lem} 


 


Действительно, 


\begin{align*} 


M(\sigma,F)&\le\sum^\infty_{n=1}|b_n|\exp(\sigma\lambda_n)= 


\sum^\infty_{n=1}|b_n|^{1-q}\biggl\{|b_n|\exp\biggr( 


\frac{\sigma+p}{q}\lambda_n\biggr)\biggr\}^q\exp(-p\lambda_n)\le \\ 


&\le\mu\biggl(\frac{\sigma+p}{q}\biggr)^q\sum^\infty_{n=1} 


|b_n|^{1-q}\exp(-p\lambda_n)=K_0\mu 


\biggl(\frac{\sigma+p}{q},\,F \biggr)^q.


\end{align*} 


 


Нам будут нужны также следующие леммы. 


 


\begin{lem}[{\series{m}\shape{n}\selectfont{}[3]}] Из каждой 


положительной возрастающей последовательности 


$\mu_n$ такой, что 


$$ 


\underset{n\to\infty}{\varlimsup}\frac{1}{\mu_n} 


\ln n\ge a>0, 


$$ 


можно выделить подпоследовательность $(\mu^*_k)$ такую, что 


$\ln k\le a\mu^*_k+1$ 


для всех $k\in \Bbb N$ и $\ln k_j\ge a\mu^*_{k_j}$ для некоторой 


возрастающей 


последовательности $(k_j)$ натуральных чисел. 


\end{lem} 


 


\begin{lem}[{\series{m}\shape{n}\selectfont{}[4], с.20--22}] Пусть 


$\varphi$ --- положительная дважды непрерывно 


дифференцируемая на $[a,+\infty)$ функция такая, что 


$\varphi'(x)\to+\infty$ 


$(x\to+\infty)$, $x^2\varphi''(x)\to+\infty$ $(x\to+\infty)$,


$\varphi''(u(x))\sim\varphi''(x)$ при $x\to+\infty$,


$|u(x)-x|\le\frac{A}{\sqrt{p''(x)}}$, где $A$ --- произвольное 


положительное 


число. Тогда функция $H(x;\lambda)=x\lambda-\varphi(x)$ 


имеет при больших значениях 


$\lambda$ единственную точку максимума $x(\lambda)$ и 


\begin{equation} 


\int^\infty_a e^{H(x;\lambda)}dx=(1+o(1)) 


\sqrt{\frac{2\pi}{p''(x(\lambda))}}\, 


e^{H(x(\lambda);\lambda)},\quad \lambda\to+\infty. 


\end{equation} 


\end{lem} 


 


{\bf 2.} Пусть сначала $0<A<+\infty$. Через $\Pi_D(A,h)$ 


обозначим класс 


функций $F\in \Pi_D(A)$ таких, что 


$$ 


\underset{n\to\infty}{\varlimsup} 


\frac{\ln n}{\ln(|b_n|\exp\{A\lambda_n\})}\le h 


\in[0,+\infty). 


$$ 


 


\begin{thm} Пусть $\beta\in[0,+\infty)$. 


Для того чтобы для любой функции 


$F\in \Pi_D(A,h)$ выполнялось неравенство 


\begin{equation} 


M(\sigma,F)\le K_0\mu 


\biggl(\frac{\sigma+A\beta}{1+\beta},\,F \biggr)^{1+\beta}, 


\quad K_0=\const >0, 


\end{equation} 


для всех $\sigma\in[0,A)$, необходимо и достаточно, чтобы $\beta>h$. 


\end{thm} 


 


\begin{pf} Пусть $\beta>h$ и $F$ --- произвольная функция из 


$\Pi_D(A,h)$. Положим $p=\beta A$ и $q=1+\beta$. Тогда 


$$ 


\underset{n\to\infty}{\varliminf} 


\frac{(q-1)\ln|b_n|+p\lambda_n}{\ln n}= 


\underset{n\to\infty}{\varliminf} 


\frac{\beta(\ln|b_n|+A\lambda_n)}{\ln n} 


\ge\frac{\beta}{h}>1, 


$$ 


откуда легко следует выполнение условия (4) леммы 2. Так как 


$$ 


0<\frac{\sigma+p}{q}=\frac{\sigma+\beta A}{1+\beta}<A 


$$ 


для всех $\sigma\in[0,A)$, то по лемме 2 имеет место неравенство (5), 


которое 


в данном случае равносильно неравенству (7). Достаточность 


условия $\beta>h$ доказана. 


 


Чтобы доказать его необходимость, нужно показать, что существует 


функция $F_1\in\Pi_D(A,h)$ такая, что для любого $\beta\le h$ 


неравенство (7) не выполняется. 


 


Пусть сначала $h>0$, а 


\begin{equation} 


F_1(s)=\sum^\infty_{n=2}n^{1/h} 


e^{-A\ln n\ln\ln n}e^{s\ln n\ln\ln n}. 


\end{equation} 


Так как здесь $b_n=n^{1/h}\exp\{-A\ln n\ln\ln n\}$, то легко видеть, что 


$\tau=0$ 


и $\gamma=1/A$. Поэтому по лемме 1 абсцисса абсолютной сходимости 


ряда (8) равна $A$, и легко видеть, что $F_1\in\Pi_D(A,h)$. 


 


Обозначим $\Phi(t)=\frac{1}{h}t-(A-\sigma)t\ln t$, 


$0\le\sigma<A$. Нетрудно 


проверить, 


что функция $\Phi$ вогнутая на $[0,+\infty)$, имеет единственную точку 


максимума 


$t(\sigma)=\frac{1}{e}\exp\Bigl\{\frac{1}{h(A-\sigma)}\Bigr\}$ и 


$\Phi(t(\sigma))=\frac{A-\sigma}{e}\exp\Bigl\{\frac{1}{h(A-\sigma)} 


\Bigr\}$. 


Отсюда 


следует, во-первых, что 


\begin{multline} 


\mu(\sigma,F_1)=\max\{e^{\Phi(\ln n)}:n\ge 2\} 


\le\exp(\max\{\Phi(t):t\ge 0\})=\\ 


=\exp\biggl(\frac{A-\sigma}{e}\exp\biggl\{ 


\frac{1}{h(A-\sigma)}\biggr\}\biggr),\quad 


0\le\sigma<A, 


\end{multline} 


и, во-вторых, 


\begin{multline} 


M(\sigma,F_1)=F_1(\sigma)=\sum^\infty_{n=2}\exp\{\Phi(\ln n)\}\ge 


\int^\infty_1\exp\{\Phi(\ln t)\}dt-\exp\Phi(t(\sigma))=\\ 


=\int^\infty_0\exp\{\Phi(t)+t\}dt-\exp 


\biggl(\frac{A-\sigma}{e}\exp\biggl\{\frac{1}{h(A-\sigma)}\biggr\} 


\biggr). 


\end{multline} 


Произведем в последнем интеграле замену $t=\frac{x}{A-\sigma}$. Тогда 


$$ 


\int^\infty_0\exp\{\Phi(t)+t\}dt=\frac{1}{A-\sigma} 


\int^\infty_0\exp\biggl\{ 


\biggl(\frac{1+h}{h(A-\sigma)}-\ln\frac{1}{A-\sigma}\biggr) 


x-x\ln x\biggr\}dx 


$$ 


и, если положим 


$\lambda=\frac{1+h}{h(A-\sigma)}-\ln\frac{1}{A-\sigma}$ и 


$\varphi(x)=x\ln x$,


то функция $\varphi$ 


удовлетворяет условиям леммы~4, 


$x(\lambda)=\frac{1}{\varphi''(x(\sigma))}=H(x(\lambda);\lambda)= 


e^{\lambda-1}$. 


Поэтому 


\begin{multline*} 


\int^\infty_0\exp\{\Phi(t)+t\}dt= 


\frac{1+o(1)}{A-\sigma}\sqrt{\frac{2\pi}{e}e^\lambda}\, 


\exp\biggl\{\frac{1}{e}e^\lambda\biggr\}=\\ 


=(1+o(1))\sqrt{\frac{2\pi}{e(A-\sigma)}\exp\frac{1+h}{h(A-\sigma)}}\; 


\exp\biggl(\frac{A-\sigma}{e}\exp\frac{1+h}{h(A-\sigma)} \biggr) 


\end{multline*} 


при $\sigma\to A$. Поэтому из (9) и (10) при $\beta\le h$ имеем 


\begin{multline*} 


\frac{M(\sigma,F_1)}{\mu\Bigl( 


\frac{\sigma+\beta A}{1+\beta},F_1\Bigr)^{1+\beta}}\ge (1+o(1)) 


\frac{\sqrt{\frac{2\pi}{e(A-\sigma)}\exp 


\frac{1+h}{h(A-\sigma)}}\exp\Bigl(\frac{A-\sigma}{e} 


\exp\frac{1+h}{h(A-\sigma)}\Bigr)} 


{\exp\Bigl(\frac{A-\sigma}{e}\exp\frac{1+h}{h(A-\sigma)}\Bigr)}\ge \\ 


\ge(1+o(1)) 


\sqrt{\frac{2\pi}{e(A-\sigma)}\exp\frac{1+h}{h(A-\sigma)}}\; 


\longrightarrow +\infty \quad (\sigma\to A), 


\end{multline*} 


т.е. неравенство (7) не выполняется. 


 


Используя лемму 4, можно также показать, что для принадлежащей 


классу $\Pi_D(A,0)$ функции 


$$ 


%% tts ??                vvv


F_1(\sigma)=\sum^\infty_{n=9} 


e^{\ln n\ln\ln\ln n} 


e^{-A\ln n\ln\ln n}e^{\sigma\ln n\ln \ln n} 


$$ 


неравенство (7) с $\beta=0$ не выполняется. 


\end{pf} 


 


{\bf 3.} Перейдем к рассмотрению случая $A=+\infty$. Через 


$\Pi_D(\infty,h)$ 


обозначим класс функций $F\in\Pi_D(\infty)$ таких, что 


$$ 


\underset{n\to\infty}{\varliminf} 


\frac{\ln n}{\ln(1/|b_n|)}\le h\in[0,1). 


$$ 


 


\begin{thm} Пусть $\beta\in[0,1)$. Для того чтобы для любой функции 


$F\in\Pi_D(\infty,h)$ выполнялось неравенство 


\begin{equation} 


M(\sigma,F)\le K_0\mu\biggl(\frac{\sigma}{1-\beta},\,F 


\biggr)^{1-\beta},\quad K_0=\const >0,


\end{equation} 


для всех $\sigma\ge 0$, необходимо и достаточно, чтобы $\beta>h$. 


\end{thm} 


 


\begin{pf} Пусть $h<\beta<1$ и $F\in\Pi_G(\infty,h)$. Выберем $p=0$ и 


$q=1-\beta$. Тогда 


$$ 


\underset{n\to\infty}{\varliminf} 


\frac{(q-1)\ln|b_n|+p\lambda_n}{\ln n}= 


\underset{n\to\infty}{\varliminf} 


\frac{\beta\ln 1/|b_n|}{\ln n}\ge\frac{\beta}{h}>1, 


$$ 


т.е. выполнено условие (4) леммы 2, и поскольку 


$\frac{\sigma}{1-\beta}\ge 0$ при 


$\sigma\ge 0$, то для всех $\sigma\ge 0$ имеет место неравенство (3), 


которое в данном 


случае равносильно неравенству (11). Достаточность условия $\beta>h$ 


доказана. 


 


Пусть теперь $0<h<1$. Рассмотрим функцию 


\begin{equation} 


F_2(\sigma)=\sum^\infty_{n=2}n^{-1/h}e^{\sigma\ln\ln n}. 


\end{equation} 


Так как здесь $\tau=h\in(0,1)$ и $\gamma=0$, то по лемме 1 абсцисса 


абсолютной сходимости ряда (12) равна $+\infty$, и легко видеть, что 


$F_2\in\Pi_D(\infty,h)$. 


 


Обозначим $\Phi(t)=-\frac{1}{h}e^t+\sigma t$, $\sigma\ge 0$. Функция 


$\Phi$ вогнутая на 


$(-\infty,+\infty)$, имеет единственную точку максимума 


$t(\sigma)=\ln(h\sigma)$ и 


$\Phi(t(\sigma))=\sigma\ln\frac{\sigma h}{e}$. Поэтому, как при 


доказательстве теоремы 1, имеем 


\begin{equation} 


\mu(\sigma,F_2)\le \frac{\sigma h}{e},\quad 


\sigma\ge 0, 


\end{equation} 


и 


\begin{equation} 


M(\sigma,F_2)+\biggl(\frac{\sigma n}{e} \biggr)^\sigma \ge 


\int^\infty_2 e^{\Phi(\ln\ln t)}dt= 


\int^\infty_{\ln\ln 2}e^{\Phi(t)}e^{e^t}e^t\,dt\ge 


\int^\infty_0\exp\biggl\{(\sigma+1)s- 


\frac{1-h}{h}e^x\biggr\}dx. 


\end{equation} 


 


Полагая $\lambda=\sigma+1$ и $\varphi(t)=\frac{1-h}{h}e^t$, 


видим, что 


функция $\varphi$ удовлетворяет 


условиям леммы~4, $x(\lambda)=\ln\frac{h(\sigma+1)}{1-h}$, 


$\varphi''(x(\lambda))=\sigma+1$ и 


$H(x(\lambda);\lambda)=(\sigma+1)\ln\frac{h(\sigma+1)}{1-h}$. Поэтому 


$$ 


M(\sigma,F_2)\ge (1+o(1))\sqrt{\frac{2\pi}{\sigma+1}}\, 


\biggl(\frac{h(\sigma+1)}{e(1-h)} \biggr)^{\sigma+1},\quad 


\sigma\to+\infty,


$$ 


и в силу (13) при $\beta\le h$ выполняется 


\begin{multline*} 


\frac{M(\sigma,F_2)}{\mu\Bigl(\frac{\sigma}{1-\beta}, 


F_2\Bigr)^{1-\beta}}\ge 


\frac{(1+o(1))\frac{h\sqrt{2\pi(\sigma+1)}}{e(1-h)} 


\Bigl(\frac{h(\sigma+1)}{e(1-h)}\Bigr)^\sigma} 


{\Bigl(\frac{\sigma h}{e(1-\beta)}\Bigr)^\sigma}\ge \\ 


\ge (1+o(1))\frac{h\sqrt{2\pi}}{e(1-h)} 


\biggl(\frac{1-\beta}{1-h} \biggr)^\sigma\sqrt{\sigma}\to 


+\infty\quad (\sigma\to+\infty), 


\end{multline*} 


т.е. для функции (12) неравенство (11) с $\beta\le h$ не выполняется. 


 


Очевидно, функция 


$$ 


\exp e^s=\sum^\infty_{n=0}\frac{1}{n!}e^{sn} 


$$ 


принадлежит классу $\Pi_D(\infty,0)$, а для нее неравенство (11) с 


$\beta=0$ не 


выполняется. 


\end{pf} 


 


Зафиксируем теперь последовательность $B=(b_n)$ комплексных 


чисел и допустим, что 


$$ 


\underset{n\to\infty}{\varliminf} 


\frac{\ln n}{\ln(1/|b_n|)}=h(B)\in[0,1), 


$$ 


а через $\Pi^*_D(\infty,B)$ обозначим класс абсолютно сходящихся 


в $\{s:0\le\re s<\infty\}$ 


рядов Дирихле (1), показатели $\lambda_n$ которых удовлетворяют 


условиям, наложенным в начале п.1. 


 


\begin{thm} Пусть $\delta\in[0,1)$. Для того чтобы для каждой функции 


$F\in\Pi^*_D(\infty,B)$ выполнялось соотношение 


\begin{equation} 


M(\sigma,F)\le\mu 


\biggl(\frac{1+o(1)}{1+\delta}\sigma,F \biggr),\quad 


\sigma\to+\infty, 


\end{equation} 


необходимо и достаточно, чтобы $\delta\ge h(B)$. 


\end{thm} 


 


\begin{pf} Если $h(B)\le\delta$, то по теореме 2 для каждого 


$\beta\in(\delta,1)$ 


и всех $\sigma\ge 0$ имеет место неравенство (11), откуда следует


$$ 


\ln M(\sigma,F)\le\ln\mu 


\biggl(\frac{\sigma}{1-\beta},F \biggr)+\ln K_0, 


$$ 


и в силу выпуклости функции $\ln\mu(\sigma,F)$ (или $M(\sigma,F)$) 


легко получаем 


соотношение (15) с $\beta$ вместо $\delta$. Так как $\beta$ --- 


произвольное число, 


то тем самым достаточность условия $\delta\ge h(B)$ доказана. 


 


Для доказательства необходимости обозначим 


$B_n=\ln(1/|b_n|)$ и 


предположим, что $h(B)>\delta$. Тогда $B_n\uparrow+\infty$ 


($n\to\infty$) и 


$$ 


\underset{n\to\infty}{\varlimsup}\frac{\ln n}{B_n} 


=h\in(0,1),\quad h=h(B). 


$$ 


Поэтому по лемме 3 существует подпоследовательность $(B^*_k)$ такая, 


что $B^*_k\ge e^2$, 


\begin{equation} 


\ln k\le h B^*_k+1 


\end{equation} 


для всех $k\in \Bbb R$ и 


\begin{equation} 


\ln k_j\ge h B^*_{k_j} 


\end{equation} 


для некоторой возрастающей последовательности $(k_j)$ натуральных 


чисел. 


 


Положим $\lambda_n=-\infty$, если $B_n\ne B^*_k$, и 


$\lambda_n=\lambda^*_k$, если $B_n=B^*_k$, где 


$\lambda^*_k=\frac{B^*_k}{\ln\ln B^*_k}$. 


Таким образом, приходим к ряду Дирихле 


$$ 


F_3(s)=\sum^\infty_{k=1}\exp\biggl\{-B_k+s 


\frac{B_k}{\ln\ln B_k}\biggr\}, 


$$ 


где для простоты $B_k=B^*_k$. В силу (16) этот ряд целый. 


Ясно, что 


$$ 


\ln\mu(\sigma,F_3)\le\max\biggl\{t 


\biggl(\frac{\sigma}{\ln\ln t}-1 \biggr):t\ge e^2\biggr\}. 


$$ 


Этот максимум достигается в точке $t=t(\sigma)$, удовлетворяющей 


уравнению 


\begin{equation} 


\sigma(\ln t\ln\ln t-1)=\ln t\ln^2\ln t, 


\end{equation} 


а 


\begin{equation} 


\ln\mu(\sigma,F_3)\le 


\frac{\sigma t(\sigma)}{\ln t(\sigma)\ln^2\ln t(\sigma)}= 


\frac{t(\sigma)}{\ln t(\sigma)\ln\ln t(\sigma)-1}. 


\end{equation} 


Решение уравнения (18) будем искать в виде 


$t(\sigma)=\exp\{e^\sigma-\gamma(\sigma)\}$, 


где, естественно, $\gamma(\sigma)=o(e^\sigma)$, $\sigma\to+\infty$. 


Подставляя $t(\sigma)$ в (18), 


имеем 


$$ 


\sigma-\frac{\sigma}{(e^\sigma-\gamma(\sigma)) 


\ln(e^\sigma-\gamma(\sigma))}=\sigma+\ln 


\biggl(1-\frac{\gamma(\sigma)}{e^\sigma} \biggr), 


$$ 


т.е. 


$$ 


\gamma(\sigma)=(1+o(1)) 


\frac{e^\sigma \sigma}{(e^\sigma-\gamma(\sigma))\ln(e^\sigma- 


\gamma(\sigma))}=(1+o(1)),\quad \sigma\to+\infty. 


$$ 


 


Поэтому $t(\sigma)=\frac{1+o(1)}{e}\exp\{e^\sigma\}$, 


$\sigma\to+\infty$, 


и из (19) имеем 


\begin{equation} 


\ln\mu(\sigma,F_3)\le 


\frac{(1+o(1))}{e\sigma}e^{-\sigma}\exp\{e\sigma\} 


\end{equation} 


для всех $\sigma\ge\sigma_0$. Далее, используя (16) и (17), получаем 


\begin{multline*} 


M(\sigma,F_3)\ge \sum_{[k_j/2]\le k\le k_j} 


\exp\biggl\{-B_k+\sigma\frac{B_k}{\ln\ln B_k}\biggr\}\ge 


\frac{k_j}{2}\exp\biggl\{-B_{k_j}+ 


\frac{\sigma}{\ln\ln B_{k_j}}B_{[k_j/2]}\biggr\}\ge\\ 


%


\ge \exp\biggl\{h B_{k_j}-\ln 2-B_{k_j}+ 


\frac{\sigma}{\ln\ln B_{k_j}} 


\frac{\ln[k_j/2]-1}{h}\biggr\} \ge\\ 


%


\ge \exp\biggl\{(h-1)B_{k_j}-\ln 2+ 


\frac{\sigma}{\ln\ln B_{k_j}} 


\biggl(\frac{\ln k_j}{h}-\frac{3}{h} \biggr)\biggr\}\ge\\ 


\ge\exp\biggl\{(h-1)B_{k_j}+ 


%


\frac{\sigma}{\ln\ln B_{k_j}}B_{k_j}-\frac{3}{h} 


\frac{\sigma}{\ln\ln B_{k_j}}-\ln 2\biggr\} 


\end{multline*} 


для всех $\sigma\ge 0$. Взяв здесь 


$\sigma=\sigma_j=(1-h+\varepsilon)\ln\ln B_{k_j}$, где $\varepsilon$ --- 


произвольное положительное число, в силу (20) имеем 


\begin{multline*} 


M(\sigma,F_3)\ge \exp\biggl\{\varepsilon B_{k_j}- 


\frac{3(1-h+\varepsilon)}{h}-\ln 2\biggr\}=\\ 


=\exp\biggl\{\varepsilon\exp\exp\frac{\sigma_j}{1-h+\varepsilon} 


-\frac{3(1-h+\varepsilon)}{h}-\ln 2\biggr\}\ge\\ 


\ge\exp\biggl\{\frac{\varepsilon\sigma_j}{1-h+\varepsilon} 


e^{\frac{\sigma_j}{1-h+\varepsilon}}\ln\mu 


\biggl(\frac{\sigma_j}{1-h+\varepsilon},F_3 \biggr)- 


\frac{3(1-h+\varepsilon)}{h}-\ln 2\biggr\}\ge 


\mu\biggl(\frac{\sigma_j}{1-h+\varepsilon},F_3 \biggr) 


\end{multline*} 


для всех $j\ge j_0(\varepsilon)$. 


В силу произвольности $\varepsilon$ отсюда следует, 


что неравенство (15) для функции $F_3$ с $\delta<H(B)$ выполняться не 


может. 


\end{pf} 


 


\begin{thebibliography}{9} 


 


\bibitem{St} 


Стрелиц Ш.И. {\em Асимптотические свойства аналитических решений 


дифференциальных уравнений}. -- Вильнюс: Минтис, 1972. -- 468~с. 


\bibitem{S1} 


Шеремета М.М. {\em Цiлi ряди Дiрiхле}. -- 1993. -- 168 с. 


\bibitem{S2} 


Шеремета М.Н. {\em 0 поведении максимума модуля целого ряда Дирихле 


вне исключительного множества} // Матем. заметки. -- 1995. 


-- Т.\,57. -- \No\,2. -- С.\,283--296. 


\bibitem{Yev} 


Евграфов М.А. {\em Асимптотические оценки и целые функции}. -- 2-е изд. 


-- M.: Физматгиз, 1962.-- 200~с. 


 


 


\end{thebibliography} 


 


\vskip 2cm 


 


\post{\it Львовский государственный }{\it Поступила} 


\post{\it университет $($Украина$)$}{15.06.1995} 


 


\end{document} 


